\documentclass[11pt,a4paper]{article}

\usepackage[utf8]{inputenc}
\usepackage[T1]{fontenc}
\usepackage{amsmath,amssymb,amsthm}
\usepackage{booktabs}
\usepackage{hyperref}
\usepackage[margin=1in]{geometry}

\newtheorem{theorem}{Theorem}
\newtheorem{proposition}{Proposition}
\newtheorem{remark}{Remark}

\title{Improved Lower Bounds for the Hypergraph Ramsey Numbers $R(4,4,4;3)$ and $R(6,6;5)$}
\author{[Your Name]\\
\small [Your Affiliation / Independent Researcher]\\
\small \texttt{[your.email@example.com]}}
\date{November 2026}

\begin{document}
\maketitle

\begin{abstract}
We establish two improved explicit lower bounds for small hypergraph Ramsey numbers:
\[
R(4,4,4;3) \ge 102 \qquad \text{and} \qquad R(6,6;5) \ge 87.
\]
These improve upon the previous best lower bounds $R(4,4,4;3) \ge 84$ (Yang, September 2026) and $R(6,6;5) \ge 84$ (Lysenst\o en, November 2026). Both constructions are compactly specified as group-invariant colorings on $1{,}162$ triple orbits and $5{,}159$ five-subset orbits, respectively, discovered using multi-orbit stabilizer-free affine group actions, multi-apex extensions, and multi-threaded dynamic-weighting ProbSAT with iterative warm-start climbing. Each witness has been independently certified by a standalone, zero-pruning exhaustive verifier over all $\binom{101}{4} = 4{,}082{,}925$ quadruples and all $\binom{86}{6} = 470{,}155{,}077$ six-subsets, respectively.
\end{abstract}

\section{Introduction}
For positive integers $t \ge 2$ and $k_1, \dots, k_r \ge t$, the hypergraph Ramsey number $R(k_1, \dots, k_r; t)$ denotes the least integer $N$ such that every $r$-coloring of the $t$-element subsets of an $N$-element set contains a subset of size $k_i$ all of whose $t$-element subsets receive color $i$, for some $1 \le i \le r$. Explicit lower bounds and known values are maintained in Radziszowski's dynamic survey \emph{Small Ramsey Numbers}~\cite{RadziszowskiDS1}.

Recent work has achieved substantial progress on Section~7.1 of~\cite{RadziszowskiDS1} by restricting edge colorings to be invariant under transitive affine or projective groups and solving the resulting compressed orbit constraints with CDCL SAT solvers. In September 2026, Yang~\cite{Yang2026} proved $R(4,4,4;3) \ge 84$ via an index-$2$ affine group over $\mathbb{Z}_{83}$ (27 triple orbits), improving the survey bound of $79$. In November 2026, Lysenst\o en~\cite{Lysenstoen2026} improved six small two-color hypergraph Ramsey numbers, including $R(6,6;5) \ge 84$ via the full affine group $\mathrm{AGL}(1,\mathbb{Z}_{83})$ (improving the survey bound of $72$).

A practical bottleneck of CDCL-based orbit searches is that random-like symmetric $k$-SAT formulas become intractable once the number of orbit variables grows beyond several dozen to a few hundred. This forces the use of large single-orbit groups, which frequently introduce stabilizer obstructions (such as elements of order $4$ or $6$ acting transitively on a $4$- or $6$-subset, forcing monochromatic cliques) or collapse pair orbits when fixed apex vertices are adjoined.

In this note, we show that combining \emph{medium-orbit stabilizer-free group actions} (specifically, diagonal multi-orbit odd-order affine actions over prime-power fields and multi-apex extensions) with \emph{dynamic-weighting ProbSAT} and \emph{iterative apex warm-start climbing} scales smoothly to $1{,}000\text{--}5{,}500$ orbit variables on a single consumer workstation, yielding two new lower bounds.

\begin{theorem}\label{thm:main}
The following hypergraph Ramsey lower bounds hold:
\begin{enumerate}
    \item $R(4,4,4;3) \ge 102$,
    \item $R(6,6;5) \ge 87$.
\end{enumerate}
\end{theorem}

\begin{table}[h]
\centering
\begin{tabular}{lcccc}
\toprule
Ramsey Number & DS1 Rev.~18~\cite{RadziszowskiDS1} & Late 2026 Best & New Bound & Improvement \\
\midrule
$R(4,4,4;3)$ & $\ge 79$ & $\ge 84$~\cite{Yang2026} & $\mathbf{\ge 102}$ & $+18$ \\
$R(6,6;5)$   & $\ge 72$ & $\ge 84$~\cite{Lysenstoen2026} & $\mathbf{\ge 87}$ & $+3$ \\
\bottomrule
\end{tabular}
\caption{Summary of improved lower bounds relative to Radziszowski's Dynamic Survey (Rev.~18, April 2026) and late-2026 preprints.}
\label{tab:summary}
\end{table}

\section{Constructions and Methodology}

Let a finite permutation group $G$ act on a vertex set $V$ of size $n$. Any $G$-invariant $r$-coloring of $\binom{V}{t}$ assigns a single color in $\{0, \dots, r-1\}$ to each $G$-orbit of $t$-subsets. For Ramsey parameters with $k_i = t+1$ for all $i$, every $(t+1)$-subset containing a canonical representative of a $t$-orbit yields a monochromatic-avoidance hyperedge over at most $t+1$ orbit variables.

\subsection{Construction for $R(4,4,4;3) \ge 102$}
To avoid elements of order $4$ or Klein-$4$ subgroups that force monochromatic $K_4^{(3)}$ configurations, we strip the entire $2$-Sylow subgroup from the multiplicative group of a prime-power finite field $\mathbb{F}_q$ ($q$ odd). Furthermore, to permit asymmetric cross-layer clique breaking and prevent pair-orbit collapse under fixed apex vertices, we act diagonally on multiple copies of $\mathbb{F}_q$ without layer-swap symmetry.

Specifically, take $q = 49 = 7^2$ and let $H = \mathrm{AGL}^{(16)}(1, \mathbb{F}_{49}) = \{x \mapsto a x + b : a \in \langle \alpha^{16} \rangle,\; b \in \mathbb{F}_{49}\}$, where $\alpha$ is a primitive element of $\mathbb{F}_{49}^\times$. Since $|\mathbb{F}_{49}^\times| = 48$, the multiplicative subgroup $\langle \alpha^{16} \rangle$ has odd order $3$, so $|H| = 49 \times 3 = 147$ is odd (containing no involutions or order-$4$ elements), and unordered pairs in $\mathbb{F}_{49}$ split into $16$ orbits.

We let $H$ act on the $101$-element set
\[
V = (\mathbb{F}_{49} \times \{0, 1\}) \cup \{\infty_1, \infty_2, \infty_3\},
\]
acting diagonally on the two copies of $\mathbb{F}_{49}$ and fixing $\infty_1, \infty_2, \infty_3$ pointwise. The group $H$ is generated by $3$ explicit permutations (two additive basis translations of $\mathbb{F}_{49} \cong \mathbb{F}_7^2$ and multiplication by $\alpha^{16}$). Under this action:
\begin{itemize}
    \item The $\binom{101}{3} = 166{,}650$ triples partition into exactly $1{,}162$ $H$-orbits.
    \item Anchoring at canonical triple-orbit representatives produces exactly $27{,}885$ unique orbit-level $K_4^{(3)}$ constraints (constraint-to-variable ratio $24.00$).
\end{itemize}
Using 32-threaded dynamic clause-weighting ProbSAT (PAWS, break base $c_b = 3.80$, weight-bump interval $500$), a zero-violation $3$-coloring was found in $2.03$ seconds ($775{,}784$ flips). Across all $166{,}650$ triples, the color class sizes are $56{,}645$, $49{,}539$, and $60{,}466$.

\subsection{Construction for $R(6,6;5) \ge 87$}
For $t=5$ and $k=6$, avoiding transitive actions on $6$-subsets requires avoiding subgroups of order $3$ or $6$ in the affine action; hence $q = 83$ ($83 - 1 = 82 = 2 \cdot 41$, $3 \nmid 83 \cdot 82$) is a natural base field, as used in~\cite{Lysenstoen2026}.

We extend the full affine group $G = \mathrm{AGL}(1, \mathbb{F}_{83})$ (order $83 \times 82 = 6{,}806$, generated by $x \mapsto x+1$ and $x \mapsto 2x$ modulo $83$) by adjoining $r = 3$ fixed apex vertices:
\[
V = \mathbb{F}_{83} \cup \{\infty_1, \infty_2, \infty_3\}, \qquad |V| = 86.
\]
Under this action:
\begin{itemize}
    \item The $\binom{86}{5} = 34{,}826{,}302$ five-subsets partition into $5{,}159$ $G$-orbits.
    \item There are $69{,}270$ unique orbit-level $K_6^{(5)}$ constraints (ratio $13.43$ hyperedges per orbit, i.e., $26.85$ clauses per variable in 6-SAT, well below the random 6-SAT threshold).
\end{itemize}
While CDCL (CaDiCaL) timed out even on the known-satisfiable $n=83$ baseline, 32-thread calibrated ProbSAT ($c_b \approx 4.85\text{--}5.05$) solved $n = 84$ ($r=1$) from random initialization in $7.99$ seconds. We then applied \emph{iterative apex warm-start climbing}: for $r = 2$ ($n = 85$) and $r = 3$ ($n = 86$), all $G$-orbits lying within the previous $(n-1)$-vertex set were initialized directly from the certified $(n-1)$-vertex coloring, newly created apex-incident orbits were randomized, and $1\%\text{--}2\%$ of mapped orbits were perturbed. This yielded certified zero-violation $2$-colorings for $n = 85$ in $37.87$ seconds and $n = 86$ in $5.72$ seconds ($3{,}493{,}357$ flips). Across all $34{,}826{,}302$ five-subsets of the $n=86$ coloring, the color class sizes are $17{,}409{,}748$ and $17{,}416{,}554$.

\section{Independent Zero-Pruning Verification}
To eliminate any possibility of bug or hidden symmetry assumption in the search pipeline, both witnesses were audited by a completely standalone verification script (\texttt{verify\_witnesses.py}, included as an ancillary file alongside the JSON witnesses). The verifier imports none of the search code, checks that every generator is a bijection on $\{0, \dots, n-1\}$, reconstructs all $\binom{n}{t}$ edge colors directly from the generators and orbit table, verifies $G$-invariance on every hyperedge, and then executes a \textbf{zero-pruning exhaustive nested loop} visiting every single $k$-subset:
\begin{itemize}
    \item \textbf{$R(4,4,4;3) \ge 102$ ($n = 101$)}: All $4{,}082{,}925 / 4{,}082{,}925$ quadruples verified with $0$ monochromatic $K_4^{(3)}$ cliques, both in compiled Numba ($0.97$s) and in pure NumPy without Numba ($0.93$s).
    \item \textbf{$R(6,6;5) \ge 87$ ($n = 86$)}: All $470{,}155{,}077 / 470{,}155{,}077$ six-subsets verified via literal $6$-nested loops with $0$ monochromatic $K_6^{(5)}$ cliques ($1.53$s on an AMD Ryzen 9 7945HX).
\end{itemize}
All witnesses and verification scripts are publicly archived at \texttt{[GitHub / Zenodo DOI URL]}.

\begin{thebibliography}{9}

\bibitem{RadziszowskiDS1}
S.~P.~Radziszowski,
\emph{Small Ramsey Numbers},
Electronic Journal of Combinatorics, Dynamic Survey \#DS1, Revision~18 (April 2026).

\bibitem{Yang2026}
H.~Yang,
\emph{Three improved lower bounds for hypergraph Ramsey numbers},
preprint (September 2026).

\bibitem{Lysenstoen2026}
C.~Lysenst\o en,
\emph{New lower bounds for six small hypergraph Ramsey numbers},
arXiv:2611.25143 (November 2026).

\bibitem{BalintSchoning2012}
A.~Balint and U.~Sch\"oning,
\emph{Choosing probability distributions for stochastic local search and the role of make versus break},
in Proc.\ SAT 2012, LNCS 7317, pp.~16--29, Springer (2012).

\end{thebibliography}

\end{document}
